\section{Introducción}
\label{sec:introduccion}
% Planteo del problema (cantidades que cambian en el tiempo; la gestión de vulnerabilidades como
% sistema dinámico), justificación, objetivo general, objetivos específicos en un
% \begin{enumerate} (el verificador N2 cuenta los \item y exige que Conclusiones los responda con
% ordinales: primer, segundo, ...), y estructura del informe con \ref a cada sección.
% Tiene que ser más larga que el resumen.

[PENDIENTE]
